\scanpage{17}

\unsure{Soient|Soit} $a\in k$, et $G$ le groupe
\[
G=\left\{\left(\begin{array}{cc}
\begin{matrix}x&ay\\y&x\end{matrix}&0\\
0&\begin{matrix}u&v\\w&z\end{matrix}
\end{array}\right)\ \middle|\ 
\begin{vmatrix}x&ay\\y&x\end{vmatrix}=\begin{vmatrix}u&v\\w&z\end{vmatrix}\ne0\right\}.
\]
\[
\left(\begin{array}{cc}
\begin{matrix}x&ay\\y&x\end{matrix}&0\\
0&\begin{matrix}u&v\\w&z\end{matrix}
\end{array}\right)\rightsquigarrow
\begin{pmatrix}u&v\\w&z\end{pmatrix}/(\text{centre})\quad\text{est un hom.}
\]
sur $GP(2)$ dont le noyau est
\[
\begin{aligned}
Z&=\left\{\left(\begin{array}{cc}
\begin{matrix}x&ay\\y&x\end{matrix}&0\\
0&\begin{matrix}u&0\\0&u\end{matrix}
\end{array}\right)\ \middle|\ x^2-ay^2=u^2\ne0\right\}\\
&=\text{centre de }G=\text{radical de }G.
\end{aligned}
\]

$G$ est déf. sur $k$, \unsure{est|et} si $\text{caract.}=2$, $a\notin k^2$, $Z$ ne l'est pas.

$\exists$ exemples analogues pour chaque caract. $p\ne0$.

% archive sheet number: 16

